% !TeX document-id = {87d6827f-5870-46b7-927e-2f8eecb71b9d}
% !TeX TXS-program:compile = txs:///pdflatex/[--shell-escape]
\documentclass{beamer}
\usetheme[faculty=ped]{fibeamer}
\usepackage[utf8]{inputenc}
\usepackage{minted}
\usemintedstyle{monokai}
\usepackage{etoolbox}
\AtBeginEnvironment{minted}{\singlespacing%
	\fontsize{10}{10}\selectfont}
\usepackage[
  main=english, %% By using `czech` or `slovak` as the main locale
                %% instead of `english`, you can typeset the
                %% presentation in either Czech or Slovak,
                %% respectively.
  czech, slovak %% The additional keys allow foreign texts to be
]{babel}        %% typeset as follows:
%%
%%   \begin{otherlanguage}{czech}   ... \end{otherlanguage}
%%   \begin{otherlanguage}{slovak}  ... \end{otherlanguage}
%%
%% These macros specify information about the presentation
\title{ LaTex for Scientific Writing} %% that will be typeset on the
\subtitle{Day 2} %% title page.
\author{\href{sambaiga.github.io}{Anthony Faustine} \\sambaiga@gmail.com}


%% These additional packages are used within the document:
\usepackage{ragged2e}  % `\justifying` text
\usepackage{booktabs}  % Tables
\usepackage{tabularx}
\usepackage{tikz}      % Diagrams
\usepackage{siunitx}
\usepackage{etoolbox}
\usepackage[round]{natbib}
\preto\tabular{\shorthandoff{-}}
\usetikzlibrary{calc, shapes, backgrounds}
\usepackage{amsmath, amssymb}
\usepackage{url}       % `\url`s
\usepackage{listings}  % Code listings
\usepackage[colorinlistoftodos]{todonotes}
\let\todox\todo
\renewcommand\todo[1]{\todox[inline]{#1}}
%\setbeamercovered{highly dynamic}
\setbeamercovered{transparent}
\frenchspacing
\newcounter{saveenumi}
\newcommand{\seti}{\setcounter{saveenumi}{\value{enumi}}}
\newcommand{\conti}{\setcounter{enumi}{\value{saveenumi}}}

\resetcounteronoverlays{saveenumi}

\begin{document}
  \frame{\maketitle}

  \AtBeginSection[]{% Print an outline at the beginning of sections
    \begin{frame}<beamer>
      \frametitle{Outline}
      \tableofcontents[currentsection]
    \end{frame}}

  \begin{darkframes}
    


\section{Graphics, Figures and Tables}

\begin{frame}[fragile]{Tables}{Creating Tables}
Use the tabular environment
\begin{minted}{tex}
\begin{tabular}[position]{column alignments}
...
\end{tabular}
\end{minted}

\begin{description}
	\item[postion] is optional (vertical position): [t] (top), [c]
	(center, this is default), [b] (bottom);
	\item [column alignments]: l (left-justified), c (center justified),
	and r (right-justified);
\end{description}
\end{frame}

\begin{frame}[fragile]{Tables}{Creating Tables}
The column data is separated by \textcolor{blue}{\&},row end is marked as \mintinline{tex}{\\} and \mintinline{tex}{\hline} draw a horizontal line.\\
Consider the following simple table:\\


\begin{tabular}[c]{|l|l|}
	\hline 
	\textbf{Parameter} & \textbf{Value}\\
	\hline 
	Path loss (n) & 2.5\\
	Model & Okumura-model\\
	Cell-radius & $1km$\\
	\hline
\end{tabular}

\end{frame}

\begin{frame}[fragile]{Tables}{Creating Tables}
Latex code for previous simple table\\
	\begin{minted}{tex}

\begin{tabular}[t]{|l|l|}
\hline 
\textbf{Parameter} & \textbf{Value}\\
\hline 
Path loss (n) & 2.5\\
Model & Okumura-model\\
Cell-radius & $1km$\\
\hline
\end{tabular}
	\end{minted}
	
	
\end{frame}


\begin{frame}[fragile]{Tables}{Creating Tables}
The \alert{\texttt{booktabs }} package  improve the quality  LaTeX tables.
\begin{itemize}
	\item The horizontal rules are called with \mintinline{tex}{\toprule}, \mintinline{tex}{\midrule} and \mintinline{tex}{\bottomrule} instead of \mintinline{tex}{\hline} command.
	\item The \mintinline{tex}{\cmidrule} is used for mid-rules that span specified columns.
	\item The content of the tables is filled in the same manner as before.
	\item To use this package first you need to add this code in preamble. \mintinline{tex}{\usepackage{booktabs}}
\end{itemize}
\end{frame}

\begin{frame}[fragile]{Tables}{Creating Tables}
	\alert{Example:} 
	
	\centering
	\begin{tabular}[t]{ll}
		\toprule
		\textbf{Parameter} & \textbf{Value}\\
		\midrule
		Path loss (n) & 2.5\\
		Model & Okumura-model\\
		Cell-radius & $1km$\\
		\bottomrule
	\end{tabular}
	
\end{frame}

\begin{frame}[fragile]{Tables}{Creating Tables}
	\alert{Example:} consider the following code:\
	\begin{minted}{tex}
\centering
\begin{tabular}[t]{ll}
\toprule
\textbf{Parameter} & \textbf{Value}\\
\midrule
Path loss (n) & 2.5\\
Model & Okumura-model\\
Cell-radius & $1km$\\
\bottomrule
\end{tabular}
	\end{minted}
\end{frame}



\begin{frame}[fragile]{Tables}{Creating Tables}
To draw multicolumn table like this one:\\
\centering

\begin{tabular}{llr}
	\toprule
	\multicolumn{2}{c}{Name} \\
	\cmidrule(lr){1-2}
	First name	& Last Name	& Grade \\
	\midrule
	John 		    & Doe	  		& $7.5$ \\
	Richard 	  & Miles		  & $2$ \\
	\bottomrule
\end{tabular}
\end{frame}

\begin{frame}[fragile]{Tables}{Creating Tables}
	Use the following command: \mintinline{tex}{\multicolumn{n}{alignment}{item}}
	\begin{description}
		\item [n]: is the number of columns to be spanned.
		\item [alignemnt]: is one of the \mintinline{tex}{l, r and c}.
		\item [item]: is the content.
	\end{description}
\alert{Example:}
\begin{minted}[numbersep=5pt]{tex}

\multicolumn{2}{c}{Name} 

\end{minted}

\end{frame}

\begin{frame}[fragile]{Tables}{Creating Tables}

\begin{minted}{tex}
\begin{tabular}{llr}
\toprule
\multicolumn{2}{c}{Name} \\
\cmidrule(r){1-2}
First name & Last Name	& Grade \\
\midrule
John  & Doe	& $7.5$ \\
Richard   & Miles  & $2$ \\
\bottomrule
\end{tabular}
\end{minted}
	
\end{frame}

\begin{frame}[fragile]{Tables}{Floating Tables}
 Latex	provides the \texttt{table}  environments for typesetting floating tables.
 \begin{itemize}
 	\item A table environment is set up as follows:
 	\begin{minted}{tex}
\begin{table}
\caption{title}
\label{tab:xxx}
%Place the table here
\end{table}
 	\end{minted}
 \end{itemize}
\end{frame}

\begin{frame}[fragile]{Tables}{Floating Tables}
	\begin{description}
		\item [\mintinline{tex}{\caption}] command is optional and used to set table tittle.
		\item [\mintinline{tex}{\label}] command is also optional and is used to reference
		the table’s number.
	\end{description}
\alert{Example:} To produce the following table
\begin{table}
\caption{Simulation Parameters}
\label{tab:model_parameter}
	\begin{tabular}[t]{ll}
	\toprule
	\textbf{Parameter} & \textbf{Value}\\
	\midrule
	Path loss (n) & 2.5\\
	Model & Okumura-model\\
	Cell-radius & $1km$\\
	\bottomrule
\end{tabular}
\end{table}

\end{frame}


\begin{frame}[fragile]{Tables}{Floating Tables}
\begin{minted}{tex}
\begin{table}
\caption{Simulation Parameters}
\label{tab:model_parameter}
\begin{tabular}[t]{ll}
\toprule
\textbf{Parameter} & \textbf{Value}\\
\midrule
Path loss (n) & 2.5\\
Model & Okumura-model\\
Cell-radius & $1km$\\
\bottomrule
\end{tabular}
\end{table}
\end{minted}

\alert{Guide to Making nice Table}: \href{https://www.inf.ethz.ch/personal/markusp/teaching/guides/guide-tables.pdf}{click here}
\end{frame}

\begin{frame}[fragile]{Activity}
	\begin{center}
		{\Large \textbf{Activity 1}}
	\end{center}
	
\end{frame}


\begin{frame}[fragile]{Graphics}{Include Graphics}
The easiest way to include images in your document is to use
the \texttt{graphicx} package.\\
Load the package graphicx:\mintinline{tex}{ \usepackage{graphicx}}
\begin{itemize}
	\item The image format available depend on what you’re using to compile.
	\item If you’re compiling using pdflatex (recommended), then you
	can use jpg, png, pdf, or eps files.
	\item Place the file in the same directory as your tex file, and use the
	\mintinline{tex}{\includegraphics[key-values]{imagefile}} command.
\end{itemize}
\end{frame}

\begin{frame}[fragile]{Graphics}{Include Graphics}
	\alert{Example:} 
	\mintinline{tex}{\includegraphics[scale=0.2]{images/bulb}}
	\begin{center}
     \includegraphics[scale=0.2]{images/bulb}
	\end{center}
	
\end{frame}

\begin{frame}[fragile]{Graphics}{Include Graphics}
	The image can be scaled to a specified height and/or width as follows:
	\mintinline{tex}{\includegraphics[height=2in,width=1in]{images/bulb}}
	\begin{center}
	\includegraphics[height=1.5in,width=1in]{images/bulb}
	\end{center}
\end{frame}

\begin{frame}[fragile]{Graphics}{Floating images}
	Use \texttt{figure} environment:
	\begin{center}
		\begin{minted}{tex}
		\begin{figure}
			\includegraphics{file}
			\caption{title }
			\label{fig:xxx}
		\end{figure}
	\end{minted}
	\end{center}
\end{frame}



\begin{frame}[fragile]{Graphics}{Floating images}
		\begin{figure}
		\includegraphics[scale=0.15]{images/bulb}
		\caption{Green Bulb} 
		\label{fig:bulb}
		\end{figure}
	To print the list of figures and tables use \mintinline{tex}{ \listoffigures} and \mintinline{tex}{\listoftables}  respectively.\\		
	
\end{frame}

\begin{frame}[fragile]{Activity}
	\begin{center}
		{\Large \textbf{Activity 2}}
	\end{center}
	
\end{frame}


\begin{frame}[fragile]{Notes with todonotes}
The \mintinline{tex}{\todo} command from the \texttt{todonotes} package is great for leaving notes to yourself and your collaborators.\\

Include this package in preamble: \mintinline{tex}{\usepackage[colorinlistoftodos]{todonotes}}


\alert{Example:}\\
\mintinline{tex}{\todo{Plain todonotes.}}  \\
	

\end{frame}

\begin{frame}[fragile]{Notes with todonotes}
\begin{itemize}
	\item Only inline notes are supported with
	beamer, but margin notes are supported for normal
	documents.
	\item There is also a handy \mintinline{tex}{\listoftodos} command.
	\begin{itemize}
		\item To use this load: \mintinline{tex}{\usepackage[colorlinks]{hyperref}} before \mintinline{tex}{\usepackage[colorinlistoftodos]{todonotes}}
	\end{itemize}
\end{itemize}
	
\end{frame}

\begin{frame}[fragile]{Activity}
	\begin{center}
		{\Large \textbf{Activity 3}}
	\end{center}
	
\end{frame}

\section{Code listing}

\begin{frame}[fragile]{Code listing}
LaTeX allows formatting and highlighting source code in your documents.
\begin{itemize}
	\item Several packages available: https://www.sharelatex.com/learn/Code\_listing
	\item Minted package  is pretty awesome and more powerful than the other  packages.
	\item  Minted uses the Python library Pygments
\end{itemize}
\end{frame}


\begin{frame}[fragile]{Code listing}
	To use Minted \href{https://www.sharelatex.com/learn/Code_Highlighting_with_minted}{Minted}
	\begin{itemize}
		\item Load the minted package: \mintinline{tex}{\usepackage{minted}}
		\item In windows add the following line at the top of your tex file.
		% !TeX TXS-program:compile = txs:///pdflatex/[--shell-escape]
		\item  In-line listing: \mintinline{tex}{\mintinline{language}{code}}
		\item block listing: Use \mintinline{tex}{\begin{minted}{code}...\end{minted}} environment.
	\end{itemize}
\end{frame}

\begin{frame}[fragile]{Code listing}
	Example: Block-listing: $C++$
\begin{minted}{cpp}
#include<iostream>
using namespace std;

int main()
{
int num,factorial=1;
cout<<" Enter Number To Find Its Factorial:  ";
cin>>num;
for(int a=1;a<=num;a++)
{
factorial=factorial*a;
}
cout<<"Factorial of Given Number is ="<<factorial<<endl;
return 0;
}
\end{minted}
\end{frame}


\begin{frame}[fragile]{Code listing}
	Example: Python
	\begin{minted}{python}
     x = 2
     y = 3
print{x + y}
\end{minted}


\end{frame}

\section{Bibliography}

\begin{frame}[<+->]{Bibliography}
To manage and include references in a LATEX document use \textsc{BibTeX}.\\
\begin{description}[long text]
	\item[BibTex]: a bibliographic tool that is used with LaTeX to help organize the user's references and create a bibliography.
	\begin{itemize}
		\item A BibTeX user creates a bibliography file with \mintinline{tex}{.bib} extension.
		\item The \mintinline{tex}{.bib} file is called a BibTEX database $\Rightarrow$ a text file containing data in a structured format.
		\item Each entry in the \mintinline{tex}{.bib} file is formatted with a certain structure and is given a "key" by which the author can refer to it in the source file.
	\end{itemize}
\end{description}
\end{frame}

\begin{frame}[fragile]{Bibliography}{The BibTEX Format}
	The generic form of a BibTEX entry is
	\begin{center}
	\begin{minted}{tex}
@type{key, field1 = ‘‘ ’’ or {} or none,
 field2 = ‘‘ ’’ or {} or none, 
 ... 
 fieldn = ‘‘ ’’ or {} or none 
 }
	\end{minted}
\end{center}
\end{frame}

\begin{frame}[fragile]{Bibliography}{The BibTEX Format}
	\alert{Example:}
	\flushleft{
\begin{minted}{tex}
@article{Gettys90,
author = {Jim Gettys and Phil Karlton and Scott McGregor},
title = {The {X} Window System, Version 11},
journal = {Software Practice and Experience},
volume = {20},
number = {S2},
year = {1990},
abstract = {A technical overview of the X11 functionality.  This is an update
of the X10 TOG paper by Scheifler \& Gettys.}
}
\end{minted}
}
\end{frame}

\begin{frame}[<+->]{Bibliography}{Export .bib file from Mendeley}
\begin{itemize}
	\item Open Mendeley, and within "My Library" found on the left, select references that you would like to texport to BibTeX.
	\item In the drop-down menu in the toolbar at the top of the screen, click "File $\rightarrow$ Export.
	\item In the dropdown list of filetypes chose "BibTeX (*.bib)" and save to the same location as the LaTeX file.
\end{itemize}
 
\alert{More on Mendeley and LaTeX}: http://libguides.mit.edu/c.php?g=176186\&p=1159535\#3
\end{frame}


\begin{frame}[<+->]{Bibliography}{Auto-syncing from Mendeley to BibTeX}
	Mendeley has the built-in capability to auto sync a BibTeX file when changes have been made to your Mendeley library.\\
	To set up the Mendeley auto sync:
	\begin{itemize}
		\item Go to Mendeley Desktop preferences.
		\item Select the BibTeX tab. 
		\item Select the box labeled “Enable BibTeX syncing” and select the BiBTeX file option you prefer.
		\item Select the location where you want the generated .bib file(s) to be stored (this should be the same location as your LaTeX file(s)).
	\end{itemize}
\end{frame}

\begin{frame}[fragile]{Activity}
	\begin{center}
		{\Large \textbf{Activity 4}}
	\end{center}
	
\end{frame}

\begin{frame}[<+->]{Bibliography}{Using BibTEX in your LATEX Document}
	To use .bib file in  latex document:
	\begin{itemize}
		\item We can use \mintinline{tex}{natbib} packages with \mintinline{tex}{\citet and \citep} commands $\Rightarrow$ \href{http://merkel.texture.rocks/Latex/natbib.php}{Reference sheet for natbib usage}.
		\item Load with \mint{tex}{\usepackage[options]{natbib}} See list of  at the end of \href{http://merkel.texture.rocks/Latex/natbib.php}{Reference sheet for natbib usage}.
		\item Example: \mintinline{tex}{\usepackage[round]{natbib}}
		\item Include .bib file at the end of document with \mintinline{tex}{\bibliography{bib file}} and specify a bibliographic styles \mintinline{tex}{\bibliographystyle{stylename}}.
		
	\end{itemize}
\end{frame}


\begin{frame}[fragile]{Bibliography}{The BibTEX Format}
	\alert{Example:}
	
	\begin{minted}{tex}

	
	\bibliographystyle{apa} %\bibliographystyle{apacite} 
	\newpage
	\bibliography{bib/References_NILM}
		\end{minted}
\end{frame}

\begin{frame}[fragile]{Bibliography}{The BibTEX Format}
	\alert{Example:} 
	\begin{itemize}
		\item According to \citet{Barker2015} \ldots $\Rightarrow$ \mintinline{tex}{According to \cite{Alcala2015} \ldots}
		\item \ldots energy is important \citep{Barker2015}  $\Rightarrow$ \mintinline{tex}{\ldots energy is important \citep{Barker2015}}
	\end{itemize}	
\end{frame}

\begin{frame}[fragile]{Activity}
	\begin{center}
		{\Large \textbf{Activity 5}}
	\end{center}
	
\end{frame}


\section{Journal and Thesis}



\begin{frame}[<+->]{Thesis with Latex}
	\alert{Folder structure}: create a new folder (your project directory). 
	\begin{itemize}
		\item  Add some additional folders within this folder:
		\begin{description}
			\item [images] $\Rightarrow$ will contain all images.
			\item [tex] $\Rightarrow$ will contain \textsc{.tex} file.
			\item [bib] $\Rightarrow$ will contain bibliography files.
		\end{description}
	   \item This will help you keep the overview about your files.
	\end{itemize}
\end{frame}


\begin{frame}[fragile]{Thesis with Latex}
	\alert{Creating the main LaTeX document}: create a main tex file using document class report and save it into tex folder as Thesis.tex.
	\begin{minted}{tex}
\documentclass[a4paper,12pt]{report} 

\begin{document}

Hello World!

\end{document}
	\end{minted}
\end{frame}

\begin{frame}[fragile]{Thesis with Latex}
	\alert{Document structure}: Let’s create the main document structure for our thesis as shown below.
	\begin{columns}
		\begin{column}{0.5\textwidth}
	\begin{itemize}
		\item Title page
		\item Abstract
		\item Table of contents
		\item List of Algorithms
		\item List of Figures
		\item List of Tables
		\item Introduction
	\end{itemize}
     \end{column}
		\begin{column}{0.5\textwidth}
			\begin{itemize}
			\item Literature Review
		\item Research Methodology
		\item Discussion		
		\item Conclusion
		\item Acknowledgment
		\item Appendices
		\item Bibliography
	\end{itemize}
\end{column}
\end{columns}
\end{frame}

\begin{frame}[fragile]{Thesis with Latex}
	\begin{itemize}
		\item First create the table of contents inside the begin and end document: \mintinline{tex}{\tableofcontents} $\Rightarrow$ a headline with Content will appear.
		\item To organize your files, create a new .tex file for each chapter of your thesis.
		\begin{itemize}
			\item Lets create a new file and save it as Introduction.tex into the tex subfolder.
			\begin{minted}{tex}
\chapter{Introduction}
\lipsum[2]
\section{Motivation}
\lipsum[4]
\section{Problem Statement}
\lipsum[3]
\section{Objectives}
\lipsum[2]
			\end{minted}
		\end{itemize}
	\end{itemize}
\end{frame}

\begin{frame}[<+->]{Thesis with Latex}
	\begin{itemize}
		\item Import the file into you main document (Thesis.tex) after \mintinline{tex}{\tableofcontents} as: \mintinline{tex}{\chapter{Introduction}
\lipsum[2]
\section{Motivation}
\lipsum[4]
\section{Problem Statement}
\lipsum[3]
\section{Objectives}
\lipsum[2]} 
		\item Great, we just started creating our document structure. 
		\item Lets put all the other chapters in there as well.
		\mintinline{tex}{\chapter{Literature Review}
\lipsum[2]
}\\
		\mintinline{tex}{\chapter{Research Methodology}

\lipsum[2-4]
}\\
		\mintinline{tex}{\chapter{Results and Discussion}

\begin{algorithm}[H]
 \KwData{this text}
 \KwResult{how to write algorithm with \LaTeX2e }
 initialization\;
 \While{not at end of this document}{
  read current\;
  \eIf{understand}{
   go to next section\;
   current section becomes this one\;
   }{
   go back to the beginning of current section\;
  }
 }
 \caption{How to write algorithms}
\end{algorithm}
}\\
		\mintinline{tex}{\chapter{Conclusion and Future Works}
\lipsum[2-4]
}
		\item Don’t forget to create the files, \mintinline{tex}{\chapter{Some Name}}  and save them.
	\end{itemize}
\end{frame}

\begin{frame}[<+->]{Thesis with Latex}
	\begin{itemize}
		\item Next add the additional lists just after \mintinline{tex}{\tableofcontents}\\
		\mintinline{tex}{\listoffigures}\\
		\mintinline{tex}{\listoftables}\\
		\mintinline{tex}{\listofalgorithms}
		\item The list of algorithms needs package: \mintinline{tex}{\usepackage{algorithm2e}}
	\end{itemize}
More details on how to typeset algorithms \href{https://en.wikibooks.org/wiki/LaTeX/Algorithms}{here.} 
\end{frame}

\begin{frame}[<+->]{Thesis with Latex}
	\alert{Add Bibliography}
	\begin{itemize}
		\item Auto-sync your Mendely bibliography library to bib subfolder.
		\item Add it to your main tex file and define a bibliography style just before \mintinline{tex}{\end{document}} \\
		\mintinline{tex}{\bibliographystyle{apa}}\\
		\mintinline{tex}{\bibliography{bib/References}}
		\item Include the natbib package in preamble for citation \mintinline{tex}{\usepackage[round]{natbib}} 
		\item For example to cite, write in any of you chapter files: \mintinline{tex}{	We refer to \citet{Mvuma2016} for things you ...}
	\end{itemize}
 
\end{frame}

\begin{frame}[<+->]{Thesis with Latex}
	\alert{The title page}
	\begin{itemize}
		\item  Add a title page. Most universities require to use a predefined title page.
		\item Create a title.tex file and save it in tex subfolder. Define the title page as title page by inserting; \mintinline{tex}{\begin{titlepage}
				tittle page contents
		\end{titlepage}}
	\item Input the file just after \mintinline{tex}{\begin{document}} as follows:\\
		
\end{itemize}

\end{frame}

\begin{frame}[fragile]{Thesis with Latex}
	\alert{Sample Tittle Page}
	\begin{minted}[baselinestretch=1,fontsize=\footnotesize]{tex}
\begin{titlepage}
\centering
{\Large \textbf{ MAXIMIZING LTE PERFORMACE WITH MIMO SYSTEMS}} \\
\vspace{0.5in}
\includegraphics[width=0.4\textwidth]{images/logo} \\
\vspace{1in}
{\large \textbf {\textbf{BY}}} \\
\vspace{.5in}
{\large \textbf {\textbf{JAMES KALUMUNA}}} \\
\vspace{1in}
{\large A dissertation in (partial) fulfillment of the requirements for the degree of
Master of Science in Telecommunications Engineering of the University of Dodoma} \\
\vspace{0.5in}
{\large University of Dodoma } \\
\vspace{0.1in}
{\large Jun 2015}
\end{titlepage}

	\end{minted}
	
\end{frame}


\begin{frame}[fragile]{Thesis with Latex}
	\alert{Add certification Page}: Create a Certification.tex file and save it in tex subfolder, define it  by inserting:
	
	\begin{minted}[baselinestretch=1,fontsize=\footnotesize]{tex}
\begin{center}
{\large \textbf{CERTIFICATION}} \\
\vspace{2in}
\end{center}
The undersigned certify that they have read and hereby recommend 
for acceptance by the University of Dodoma dissertation entitled 
\textbf{Maximizing LTE performance with MIMO Systems} in fulfillment
of the requirements for the degree of \textbf{Master of Science in 
Telecommunications Engineering} of the University of Dodoma.	
	\end{minted}
	
Input the file just after \mintinline{tex}{\input{tex/tittle}}.	
\end{frame}


\begin{frame}[fragile]{Thesis with Latex}
	\alert{Add certification Page}\\
	Includes supervisors approval as follows
	
	\begin{minted}[baselinestretch=1,fontsize=\footnotesize]{tex}
	\vspace{1in}
	\centering
	Approved by: 
	\bigbreak
	
\noindent\begin{tabular}{ll}
\makebox[2.5in]{\hrulefill} & \makebox[2.5in]{\hrulefill}\\
Prof A.N Mvuma (First Supervisor) & Date\\[8ex]% adds space between the two sets of signatures
\makebox[2.5in]{\hrulefill} & \makebox[2.5in]{\hrulefill}\\
Dr. Hector Mongi (Second Supervisor) & Date\\
\end{tabular}
		
	\end{minted}	
\end{frame}

\begin{frame}[fragile]{Thesis with Latex}
Add empty pages between the title, abstract and table of contents and change the numbering to be roman.	The real Arabic page numbering of the thesis should start  with the first page of the introduction.
	\begin{itemize}
	\item First include the following commands into the preamble:
\begin{minted}[baselinestretch=1,fontsize=\footnotesize]{tex}
\setcounter{secnumdepth}{3}
 % determines up to what level the sectioning titles are numbered
\setcounter{tocdepth}{3}
\pagenumbering{roman}
\end{minted}
\item Then add \mintinline{tex}{\newcounter{rom}} just after the \mintinline{tex}{\begin{document}}.
	\end{itemize}
\end{frame}

\begin{frame}[fragile]{Thesis with Latex}
	\begin{itemize}
		\item Also add the following commands just before \mintinline{tex}{\tableofcontents}.
		\begin{minted}[baselinestretch=1,fontsize=\footnotesize]{tex}
		\addtocounter{rom}{1}\setcounter{page}{2}~
		\newpage\thispagestyle{plain}\setcounter{page}{3}
		\end{minted}
		\item Finally add the following command just before \mintinline{tex}{\section{Introduction}
Worldwide electrical energy consumption is constantly increasing especially in developing countries \citep{Monacchi2013}. According to United States Energy Information Administration report\footnote{http://www.eia.gov/todayinenergy/detail.cfm?id=14011}, \textit{the developing countries will account for 65\% of the global energy consumption by 2040}. 

\lipsum[1]

\subsection{Background}

Over the recently most part of the world has witnessed rapidly increasing to energy use in buildings (residential and commercial). Building contribute about 40\% of the total global energy \citep{Batra2014c}. Residential and commercial buildings consume approximately 60\% of the world’s electricity \footnote{The United Nation’s Environment Programme’s Sustainable Building and Climate Initiative (UNEP-SBCI)}. In U.S. A for example 74.9\% of all the electricity produced is used just to operate buildings \footnote{\href{http://www.eia.gov/todayinenergy/detail.cfm?id=14011}{United States Energy Information Administration report}} while in Africa, 56\% of the total electric power demand is from buildings \citep{Kitio2013}. Thus, energy saving in buildings will have significant impact on the reduction of overall energy consumption.

\lipsum[1]

\subsection{Problem Statement}

The key challenge to NILM problem is how to design efficient unsupervised NILM algorithm that can run in real-time using low-frequency sampling data. Several state-of-the-art NILM algorithms have been proposed using different approaches such as different variants of Hidden Markov Models (HMM) \citep{Kim2011,Parson2012,Kolter2012,Makonin2015}, Deep Neural Networks (DNN) \citep{Badayos2015,Paulo2016a}, Graph Signal Processing (GSP) \citep{Stankovic2014,Zhao2016a} and Combinatorial Optimization (CO) \citep{ReyesLua2015,Batra2013}. However, most of these algorithms suffer from high computational complexity which make them unstable for real-time applications, cannot be generalized across different buildings, requires lot of training data, their performance is limited to few numbers of appliances and are sensitive to noise and similar devices

\lipsum[1]
\subsection{Research Objectives}

\subsubsection{Main Objective}

The broad aim of this study is to develop NILM framework for sustainable residential buildings energy. The NILM framework will be well suited for the inherent characteristics of grids in Tanzania.

\subsubsection{Specific Objectives}
Specific Objectives are:
\begin{enumerate}
	\item To develop tools that will enable disaggregation research in developing countries.
	\item To develop innovative and real-time unsupervised NILM algorithms for sustainable energy consumption in residential buildings.
	\item To demonstrate and evaluate the potential of the proposed algorithm in (2) for sustainable
	energy consumption in residential buildings .
\end{enumerate}

\subsection{Research Questions}
This research is intended to answer the following questions:
\begin{enumerate}
	\item What tools can be developed to increase energy disaggregation research in developing countries?.
	\item How to design an efficient and real-time NILM algorithms that can be generalized across buildings  by taking into consideration developing countries characteristics?.
	\item Which and how innovative sustainable energy saving applications in residential buildings could be enabled with NILM algorithms in (2)?.
\end{enumerate}

\subsection{Significance of the Research}
The major contribution of the proposed study will be  novel unsupervised algorithm for energy disaggregation problem and its applicability in helping households to achieve quantifiable energy saving. The algorithm will provide real-time appliance specific information that will increase public awareness and make them be part and parcel of energy conservation. It will further help utility and policy makers gain better insights into energy consumption in residential buildings. 

The proposed study will establish tools and resource pertaining to energy consumption data sets. This will facilitate and promote research activities in energy disaggregation, energy data analysis, electricity grid modelling and appliance usage behaviour. Apart from that, energy consumption data sets  will be useful for policy makers in the energy sector.

The study will also contribute to the  understanding of challenges and possible strategies for energy conservation in residential buildings. Finally, the proposed study is expected to provide better theoretical understanding of NILM and its applicability in sustainable energy in residential buildings. }.
		\begin{minted}[baselinestretch=1,fontsize=\footnotesize]{tex}
		\newpage\thispagestyle{plain}~
		\clearpage
		\pagenumbering{arabic}
		\end{minted}
	\end{itemize}
\end{frame}

\begin{frame}[fragile]{Thesis with Latex}
	\begin{minted}[baselinestretch=1,fontsize=\footnotesize]{tex}
	documentclass[a4paper,12pt]{report} 
	
	\usepackage{lipsum}  %Generate dummy text (lorem ipsum) in your document
	\usepackage{algorithm2e}  %	
	\usepackage[round]{natbib} %for citation
	\usepackage{graphicx} % for images and figures
	
	\setcounter{secnumdepth}{3}
	\setcounter{tocdepth}{3}
	\pagenumbering{roman}
	
	\begin{document}
	\newcounter{rom}	
	\input{tex/tittle} % The tittle page
	\begin{center}
{\large \textbf{CERTIFICATION}} \\
\vspace{1in}
\end{center}
The undersigned certify that they have read and hereby recommend for acceptance by the University of Dodoma dissertation entitled \textbf{Maximizing LTE performance with MIMO Systems} in fulfillment of the requirements for the degree of \textbf{Master of Science in Telecommunications Engineering} of the University of Dodoma.	

\vspace{1in}
\begin{center}
Approved by: 
\bigbreak
\noindent\begin{tabular}{ll}
	\makebox[2.5in]{\hrulefill} & \makebox[2.5in]{\hrulefill}\\
	Prof A.N Mvuma (First Supervisor) & Date\\[8ex]% adds space between the two sets of signatures
	\makebox[2.5in]{\hrulefill} & \makebox[2.5in]{\hrulefill}\\
	Dr. Hector Mongi (Second Supervisor) & Date\\
\end{tabular}
\end{center}
\newpage
	
	\addtocounter{rom}{1}\setcounter{page}{2}~
	\newpage\thispagestyle{plain}\setcounter{page}{3}
	
	\tableofcontents % Create table of comments
	
	\end{document}
	
	\end{minted}
	
\end{frame}

\begin{frame}[fragile]{Thesis with Latex}
	\begin{minted}[baselinestretch=1,fontsize=\footnotesize]{tex}

	\listoffigures % list of figures
	\listoftables  % list of tables
	\listofalgorithms % list of algorithms
	
	\newpage\thispagestyle{plain}~
	\clearpage
	\pagenumbering{arabic}
	
	\section{Introduction}
Worldwide electrical energy consumption is constantly increasing especially in developing countries \citep{Monacchi2013}. According to United States Energy Information Administration report\footnote{http://www.eia.gov/todayinenergy/detail.cfm?id=14011}, \textit{the developing countries will account for 65\% of the global energy consumption by 2040}. 

\lipsum[1]

\subsection{Background}

Over the recently most part of the world has witnessed rapidly increasing to energy use in buildings (residential and commercial). Building contribute about 40\% of the total global energy \citep{Batra2014c}. Residential and commercial buildings consume approximately 60\% of the world’s electricity \footnote{The United Nation’s Environment Programme’s Sustainable Building and Climate Initiative (UNEP-SBCI)}. In U.S. A for example 74.9\% of all the electricity produced is used just to operate buildings \footnote{\href{http://www.eia.gov/todayinenergy/detail.cfm?id=14011}{United States Energy Information Administration report}} while in Africa, 56\% of the total electric power demand is from buildings \citep{Kitio2013}. Thus, energy saving in buildings will have significant impact on the reduction of overall energy consumption.

\lipsum[1]

\subsection{Problem Statement}

The key challenge to NILM problem is how to design efficient unsupervised NILM algorithm that can run in real-time using low-frequency sampling data. Several state-of-the-art NILM algorithms have been proposed using different approaches such as different variants of Hidden Markov Models (HMM) \citep{Kim2011,Parson2012,Kolter2012,Makonin2015}, Deep Neural Networks (DNN) \citep{Badayos2015,Paulo2016a}, Graph Signal Processing (GSP) \citep{Stankovic2014,Zhao2016a} and Combinatorial Optimization (CO) \citep{ReyesLua2015,Batra2013}. However, most of these algorithms suffer from high computational complexity which make them unstable for real-time applications, cannot be generalized across different buildings, requires lot of training data, their performance is limited to few numbers of appliances and are sensitive to noise and similar devices

\lipsum[1]
\subsection{Research Objectives}

\subsubsection{Main Objective}

The broad aim of this study is to develop NILM framework for sustainable residential buildings energy. The NILM framework will be well suited for the inherent characteristics of grids in Tanzania.

\subsubsection{Specific Objectives}
Specific Objectives are:
\begin{enumerate}
	\item To develop tools that will enable disaggregation research in developing countries.
	\item To develop innovative and real-time unsupervised NILM algorithms for sustainable energy consumption in residential buildings.
	\item To demonstrate and evaluate the potential of the proposed algorithm in (2) for sustainable
	energy consumption in residential buildings .
\end{enumerate}

\subsection{Research Questions}
This research is intended to answer the following questions:
\begin{enumerate}
	\item What tools can be developed to increase energy disaggregation research in developing countries?.
	\item How to design an efficient and real-time NILM algorithms that can be generalized across buildings  by taking into consideration developing countries characteristics?.
	\item Which and how innovative sustainable energy saving applications in residential buildings could be enabled with NILM algorithms in (2)?.
\end{enumerate}

\subsection{Significance of the Research}
The major contribution of the proposed study will be  novel unsupervised algorithm for energy disaggregation problem and its applicability in helping households to achieve quantifiable energy saving. The algorithm will provide real-time appliance specific information that will increase public awareness and make them be part and parcel of energy conservation. It will further help utility and policy makers gain better insights into energy consumption in residential buildings. 

The proposed study will establish tools and resource pertaining to energy consumption data sets. This will facilitate and promote research activities in energy disaggregation, energy data analysis, electricity grid modelling and appliance usage behaviour. Apart from that, energy consumption data sets  will be useful for policy makers in the energy sector.

The study will also contribute to the  understanding of challenges and possible strategies for energy conservation in residential buildings. Finally, the proposed study is expected to provide better theoretical understanding of NILM and its applicability in sustainable energy in residential buildings. 
	\chapter{Literature Review}
\lipsum[2]

	\chapter{Research Methodology}

\lipsum[2-4]

	\chapter{Results and Discussion}

\begin{algorithm}[H]
 \KwData{this text}
 \KwResult{how to write algorithm with \LaTeX2e }
 initialization\;
 \While{not at end of this document}{
  read current\;
  \eIf{understand}{
   go to next section\;
   current section becomes this one\;
   }{
   go back to the beginning of current section\;
  }
 }
 \caption{How to write algorithms}
\end{algorithm}

	\chapter{Conclusion and Future Works}
\lipsum[2-4]

	
	\bibliographystyle{apa}
	\bibliography{bib/References}
	
	\end{document}
	
	\end{minted}
	
\end{frame}

\begin{frame}[fragile]{Activity}
	\begin{center}
		{\Large \textbf{Activity 6}}
	\end{center}
	
\end{frame}

\section{Presentation Slides: LaTeX Beamer}

\begin{frame}[fragile]{LaTeX Beamer}
	The beamer class is a LaTeX class that allows you to create a
	beamer presentation and slides.
	\begin{block}{The basic steps to create a beamer presentation }
		\begin{itemize}
		\item Specify beamer as document class instead of article.
		\item Structure your LaTeX text using section and subsection
		commands.
		\item Place the text of the individual slides inside frame commands.
	\end{itemize}
	\end{block}

\end{frame}

\begin{frame}[fragile]{Frames}
	
Each beamer project is made up of a series of frames defined by \mintinline{tex}{\begin{frame}..\end{frame}} environment.
\begin{itemize}
	\item Each frame produces one or more slides depending on the slide overlays.
\end{itemize}
\begin{description}
	\item[The title page frame:]  simply displays a title page
	\begin{center}
            \mint{tex}{\begin{frame} \titlepage \end{frame}}
	\end{center}
    \item [The table of contents: ] dynamically creates a table of contents
    based on the sections and subsections.
 
\mint{tex}{\begin{frame}{Outline}\tableofcontents\end{frame}}

\end{description}
	     
\end{frame}


\begin{frame}[fragile]{Frame}
Basic Frame:
\mint{tex}{\begin{frame}{Tittle}
		content...
\end{frame}}

\end{frame}

\begin{frame}[fragile]{Sections and Subsections}
Presentations are divided into sections, subsections, and
sub-subsections.
\end{frame}

\begin{frame}[fragile]{Activity}
	\begin{center}
		{\Large \textbf{Activity 7}}
	\end{center}
	
\end{frame}

\begin{frame}[fragile]{Structuring a Frame}
Beamer provides many ways to structure your frames so they appear
well organized and are easy for the audience to follow.
\begin{itemize}
	\item Columns
	\item Blocks
\end{itemize}
\end{frame}

\begin{frame}[fragile]{Structuring a Frame}\framesubtitle{Coulmns}
	The column environment is called as shown below:
	\begin{center}
\begin{minted}{tex}
\begin{columns}
\column{.xx\textwidth}
First column text/graphics
\column{.xx\textwidth}
Second column text/graphics
\end{columns}
\end{minted}
\end{center}
\end{frame}

\begin{frame}[fragile]{Structuring a Frame}\framesubtitle{Coulmns}
	Here is a simple example:
		\begin{minted}{tex}
		\begin{columns}
		\column{.5\textwidth}
		Column Number 1
		\column{.5\textwidth}
		Column Number 2
		\end{columns}
		\end{minted}
\alert{Which gives us}:\\
	\begin{columns}
		\column{.5\textwidth}
		Column Number 1
		\column{.5\textwidth}
		Column Number 2
	\end{columns}	
\end{frame}

\begin{frame}[fragile]{Structuring a Frame}\framesubtitle{Blocks}
	Blocks can be used to separate a specific section of text or graphics
	from the rest of the frame.

	\begin{center}
		\begin{minted}{tex}
		\begin{block}{Block Tittle}
		content...
		\end{block}
		\end{minted}
	\end{center}
\alert{Which gives us}:\\
\begin{block}{Block Tittle}
	content...
\end{block}
\end{frame}


\begin{frame}[fragile]{Structuring a Frame}\framesubtitle{Blocks}
	Other block environments are also available:\\
\textbf{Theorem block}	
	\begin{center}
		\begin{minted}{tex}
		\begin{theorem}
		$x^2 + y^2 = 1 \rightarrow$ Circle with $r = 1$
		\end{theorem}
		\end{minted}
	\end{center}
	\alert{Which gives us}:\\
	\begin{theorem}
		$x^2 + y^2 = 1 \rightarrow$ Circle with $r = 1$
	\end{theorem}
\end{frame}

\begin{frame}[fragile]{Structuring a Frame}\framesubtitle{Blocks}
	Other block environments are also available:\\
	\textbf{Example block}	
	\begin{center}
		\begin{minted}{tex}
		\begin{example}
		This is $\ldots$
		\end{example}
		\end{minted}
	\end{center}
	\alert{Which gives us}:\\
	\begin{example}
		This is $\ldots$
	\end{example}
\end{frame}

\begin{frame}[fragile]{Structuring a Frame}\framesubtitle{Blocks}
	Other block environments are also available:\\
	\textbf{Alert block}	
	\begin{center}
		\begin{minted}{tex}
		\begin{alertblock}{Title}
		This is $\ldots$
		\end{alertblock}
		\end{minted}
	\end{center}
	\alert{Which gives us}:\\
	\begin{alertblock}{Title}
		This is $\ldots$
	\end{alertblock}
\end{frame}

\begin{frame}[fragile]{Structuring a Frame}\framesubtitle{Blocks}
	Other block environments are also available:\\
	\textbf{Lemma block}	
	\begin{center}
		\begin{minted}{tex}
		\begin{lemma}
		$x^2 + y^2 = 1$ 
		\end{lemma}
		\end{minted}
	\end{center}
	\alert{Which gives us}:\\
	\begin{lemma}
	$x^2 + y^2 = 1 $
	\end{lemma}
\end{frame}

\begin{frame}[fragile]{Structuring a Frame}\framesubtitle{Blocks}
	Other block environments are also available:\\
	\textbf{Corollary block}	
	\begin{center}
		\begin{minted}{tex}
		\begin{corollary}
		This proof
		\end{corollary}
		\end{minted}
	\end{center}
	\alert{Which gives us}:\\
	\begin{corollary}
		This proof
	\end{corollary}
\end{frame}

\begin{frame}[fragile]{Structuring a Frame}\framesubtitle{Blocks}
	Other block environments are also available:\\
	\textbf{Proof block}	
	\begin{center}
		\begin{minted}{tex}
		\begin{proof}
		This proof
		\end{proof}
		\end{minted}
	\end{center}
	\alert{Which gives us}:\\
\begin{proof}
	This proof
\end{proof}
\end{frame}

\begin{frame}[fragile]{Structuring a Frame}\framesubtitle{Columns and Blocks}

	We can combine columns and blocks to make a much cleaner looking
	presentation.
	\begin{center}
		\begin{minted}{tex}
	\begin{columns}[t]
	\column{.5\textwidth}
	\begin{block}{Column 1 Header}
	Column 1 Body Text
	\end{block}
	\column{.5\textwidth}
	\begin{block}{Column 2 Header}
	Column 2 Body Text
	\end{block}
	\end{columns}
		\end{minted}
	\end{center}

	
\end{frame}

\begin{frame}[fragile]{Structuring a Frame}\framesubtitle{Columns and Blocks}
	\alert{Which gives us}:\\
		\begin{columns}[t]
		\column{.5\textwidth}
		\begin{block}{Column 1 Header}
			Column 1 Body Text
		\end{block}
		\column{.5\textwidth}
		\begin{block}{Column 2 Header}
			Column 2 Body Text
		\end{block}
	\end{columns}


\textbf{Note}. Notice that the \mintinline{tex}{[t]} argument to the columns command top-aligned our blocks so they are vertically even as opposed to vertically centered on the slide.
\end{frame}


\begin{frame}[fragile]{Overlay}
In Beamer, overlays control the order in which parts of the frame
appear.
\begin{itemize}
	\item An easy way to implement an overlay is to place the \mintinline{tex}{\pause} command between the parts you want to show up separately.
\end{itemize}


\alert{Example}
\begin{center}
\begin{minted}{tex}
\textbf{Step1:} Compute the maximal value of $x^2$.
\pause
\textbf{Step2:} Compute the maximal value of $y^2$.
\pause
\textbf{Step3:} Add $x^2$ and $y^2$.
\end{minted}
\end{center}
\end{frame}


\begin{frame}[fragile]{Overlay}
	\alert{Which results into:}\\
	
	\textbf{Step1:} Compute the maximal value of $x^2$.\\
	\pause
	\textbf{Step2:} Compute the maximal value of $y^2$.\\
	\pause
	\textbf{Step3:} Add $x^2$ and $y^2$.
\end{frame}


\begin{frame}[fragile]{Overlay}
	\alert{Which results into:}\\
	
	\textbf{Step1:} Compute the maximal value of $x^2$.\\
	\pause
	\textbf{Step2:} Compute the maximal value of $y^2$.\\
	\pause
	\textbf{Step3:} Add $x^2$ and $y^2$.
\end{frame}

\begin{frame}[fragile]{Overlay Specification}
	Overlay specifications are given in pointed brackets (<,>) and
	indicate which slide the corresponding information should appear
	on.
	\begin{itemize}
		\item The specification \mintinline{tex}{<1->} means “display from slide 1 on.” 
		\item \mintinline{tex}{<1-3>} means “display from slide 1 to slide 3.” 
		\item  \mintinline{tex}{<2>} means “display  slide 2.” 
		\item \mintinline{tex}{<-3,5-6,8->} means “display on all slides except slides 4 and 7.”
	\end{itemize}
\textbf{Note}: If you want each item of a list to appear in order, use the \mintinline{tex}{[<+->] option}\\
E.g. \mintinline{tex}{\begin{frame}[<+->]{title} ..... \end{frame}}
\end{frame}

\begin{frame}[fragile]{Overlay } \framesubtitle{Timing Specifications With Alert}
	\alert{Consider:}
	
\begin{columns}
	\begin{column}{0.5\textwidth}
		\begin{itemize}[<+-| alert@+>]
			\item Item Aaay
			\item Item Bee
			\item Item See
			\item Item Dee
		\end{itemize}
	\end{column}
\pause
	\hspace{5pt}\vrule\hspace{5pt}%
	\begin{column}{0.5\textwidth}		
\begin{minted}{tex}
\begin{itemize}[<+-| alert@+>]
\item Item Aaay
\item Item Bee
\item Item See
\item Item Dee
\end{itemize}
\end{minted}
	\end{column}
\end{columns}		
\end{frame}




\begin{frame}[fragile]{Activity}
	\begin{center}
		{\Large \textbf{Activity 8}}
	\end{center}
	
\end{frame}

\section{Collaborative Writing}
\begin{frame}[fragile]{Collaborative writing}
Most of your writing will be collaborative:
\begin{itemize}[<+>]
	\item often participants are distributed
	\item there are lots of ways to deal with this $\Rightarrow$ even when they are local, these techniques help.
	
\end{itemize}

\end{frame}

\begin{frame}[<+>]{Collaborative writing}
	Available modes
	\begin{itemize}
		\item One person acts as editor, and incorporates changes: others  $\Rightarrow$ communicate proposed changes.
		\begin{itemize}
			\item [-] lots of work for editor, but only they end up happy.
		\end{itemize}
		\item Token: one person has the \alert{token} (for all or part)
		\begin{itemize}
			\item edit as please when have token pass it when $\Rightarrow$ finished (e.g. by email)
			\item [-] requires trust.
		\end{itemize}
		\item Truly distributed:
		\begin{itemize}
			\item [+] all have access, and can edit.
			\item [+] conflicts are merged
			\item [+-] very powerful, but requires tools.
		\end{itemize}
	\end{itemize}
\end{frame}


\begin{frame}[<+>]{Collaborative writing models
		Truly distributed collaboration}
Require tools that support
\begin{itemize}
	\item Distributed access { e.g. Dropbox} and Revision control
\end{itemize}
Available options:
\begin{description}[Example]
	\item[Payfor] \href{https://www.overleaf.com/}{Overleaf}, \href{https://www.sharelatex.com/}{ShareLaTex}, \href{https://www.authorea.com/}{Authorea} and \href{https://papeeria.com/}{Papeeria}.
	\begin{itemize}
		\item [-] often have a free limited plan.
		\item [-] focussed on latex, not the rest (e.g. accompanying code and data).
	\end{itemize}
  \item [free: standard open source tools]  Version Control eg git
\end{description}

\end{frame}

\begin{frame}[<+>]{Revision (or version) control}
Features:\begin{itemize}
	\item Allow you to see all revisions of paper $\Rightarrow$ e.g. revert back to an old version if you don't like changes.
	\item Trace activity $\Rightarrow$ volume also what changed, with comments.
	\item Good for code,  LaTeX, and (some) data.
\end{itemize}
Cloud git services: \href{https://github.com/}{Github}, \href{https://bitbucket.org/}{BitBucket},\href{https://about.gitlab.com/}{Gitlab}.
\end{frame}


\begin{frame}
	\centering
	\Huge{THANK YOU}
\end{frame}

\begin{frame}[t, allowframebreaks]

\frametitle{References}
\bibliographystyle{apa} %\bibliographystyle{apacite} 
\newpage
\bibliography{bib/References_NILM}
\end{frame} 

    \end{darkframes}


\end{document}
